Flying ad hoc networks (FANETs) require routing decisions that are both reliable under rapid topology changes and executable with only local information on resource-constrained unmanned aerial vehicles (UAVs). Centralized or communication-heavy multi-agent reinforcement learning methods can learn high-quality routing behavior, but their execution-time dependence on global state or inter-agent messages is difficult to deploy in sparse and mobile UAV swarms. This paper presents \method, a decentralized routing framework that transfers routing knowledge from a privileged global teacher to a local student and augments the student with predictive risk switching. The teacher is trained offline as a reinforcement-learning actor-critic policy with global topology information, while the deployed policy uses only one-hop candidate features, destination geometry, queue state, predicted link lifetime, onward stability, and lightweight safety gates. The proposed risk switch activates the predictive branch only when the nominal local decision is unsafe, thereby retaining the low-latency behavior of geographic routing in benign cases while improving recovery from routing holes and imminent link breaks. Full Phase 12 experiments over 14 scenarios and five random seeds show that \phaseTwelve, the fully evaluated predecessor of \method, achieves the best overall packet delivery ratio and deadline delivery among GPSR, predictive geographic routing, Evo-QGeo, IQMR Q($\lambda$), DRAMA, and internal Lite-GLOBE variants, while using fewer input bytes than predictive and message-passing RL baselines. The current Phase 13 \method{} implementation adds link-loss-aware switching, top-$k$ onward stability, energy-aware tie breaking, and drop suppression; its full-scale validation is left as an explicit experimental TODO in this draft.
